\chapter{From Worst-Case to Amortised Analysis}
\label{chap: chap6}

In this chapter, we introduce $\lambda^g_{\mathsf{amor}}$, an extension of $\lambda^g_{\mathsf{worst}}$ that supports amortised cost analysis of guarded recursive programs. 
Following $\lambda$-amor, the extended calculus incorporates an \emph{affine} type system and a \emph{graded potential} modal type. 
Together, these features let types store potential and use it to offset costs, enabling type-based amortised analysis.

We first present the extensions to the syntax, types, and semantics in Section~X.1.
In Section~X.2 we discuss the restriction to time-independent potential grades and illustrate the calculus with program examples. 
Finally, the extended model and soundness theorem of $\lambda^g_{\mathsf{amor}}$ are described in Section~X. 3.


\section{Overview of $\lambda^g_{\mathsf{amor}}$}
The formalisation of $\lambda^g_{\mathsf{amor}}$ builds upon that of $\lambda^g_{\mathsf{worst}}$ presented in Chapter\ref{chap:4}. Rather than repeating the complete specification, we introduce the new and modified constructs. 
The full formalisation is provided in Appendix~B.

\subsection{Affine Types and Potential}
Based on $\lambda$-amor, we add the potential modality $[s] \tau$, this type carries $s$ units of potential, which can be used to cancel out the cost of subsequent computations.

An important requirement is that stored potential cannot be duplicated arbitrarily, 
since this would allow the same amount of potential to offset multiple costs.
To enforce this restriction, we follow $\lambda$-amor to adopt affine types.
Specifically, variables in an affine type system are used at most once, and hence term contexts cannot be 
copied either. For example, consider the following two typing rules for function application:
\begin{mathpar}
  \inferrule*
    {\hastype{\Gamma}{e_1}{\tau_1 \to \tau_2} \quad \hastype{\Gamma,x:\tau_1}{e_2}{\tau_2}}
    { \Gamma \vdash e_1e_2 : \tau_2 }
  \and
  \inferrule*
    {\hastype{\Gamma_1}{e_1}{\tau_1 \multimap \tau_2} \quad \hastype{\Gamma_2,x:\tau_1}{e_2}{\tau_2}}
    { \Gamma_1+\Gamma_2 \vdash e_1e_2 : \tau_2 }
\end{mathpar}
The left rule is for a usual structural type system, while the right is for an affine type system, where the arrow type is written $\multimap$. Specifically, 
+ denotes the disjoint union between two contexts in the right rule, and the context must be split 
into two separate parts $\Gamma_1$ and $\Gamma_2$ to ensure no variables is used more than once.

To allow unrestricted use of variables when necessary, affine type systems employ the exponential modality $!\tau$, meaning that 
a term of type $!\tau$ can be used arbitrarily many times. This is commonly achieved by separating the typing context into an affine context and an unrestricted context.

Following this approach, the term typing judgment of $\lambda^g_{\mathsf{amor}}$ takes the form
\[
\Theta;\Delta;\Omega;\Gamma\vdash e:\tau
\]
$\Omega$ is the typing context that contains variables with unbounded usage.

\subsubsection{Terms and Typing judgment}
Figure \ref{fig:amor-typing-rules} shows selected typing rules of $\lambda^g_{\mathsf{amor}}$, which gives 
an overview of its syntactic extensions compared to $\lambda^g_{\mathsf{worst}}$.

$!e$ and $\lete{!x}{e_1}{e_2}$ are introduction and elimination constructs for the exponential type.
In particular, the introduction rule \textsc{T-ExpI} requires the affine context to be empty, ensuring that a term promoted to the exponential type does not depend on affine resources.
The elimination rule T-ExpE then allows the underlying term to be introduced into the unrestricted context, where it can be used multiple times.
 
The ordinary product type $\times$ is replaced by two types: $\otimes$ and $\&$.
The tensor pair $\ttuple{e_1}{e_2}$ typed $\tau_1 \otimes \tau_2$ can be introduced by the \textsc{T-TensorI} rule, which requires type-checking both
$e_1$ and $e_2$ by splitting the affine context. This condition is captured by its elimination form $\lete{\ttuple{x}{y}}{e_1}{e_2}$ and the \textsc{T-TensorE} rule,
where both $x$ and $y$ can be used inside the continuation $e_2$. Dually, \textsc{T-WithI} checks the ``with'' pair $\tuple{e_1}{e_2}$ typed $\tau_1\, \& \,\tau_2$ with the same affine context 
on both sides, this means using one of the components may consume all the variables in the context, and
only one side can be projected by either $\mathsf{fst}$ or $\mathsf{snd}$.

We now turn to the graded potential modality \([s]\tau\) and its associated operations. $\store e$ is the atomic operation 
that introduces potential. Potential cannot be created freely, its typing rule \textsc{E-Store} gives the result type 
$\tmonad{s}{[s]\tau}$, requiring us to first pay for the resources we want to store.
Finally, $\release{x}{e_1}{e_2}$ releases the potential stored in $e_1$ to offset the cost of the subsequent monadic computation $e_2$.
Based on its typing rule \textsc{T-Release}, $e_1$ is a pure term typed $[s] \tau$, and the stored $s$ units pay for the same amount of 
cost in $e_2$, then the final typed cost transforms from $s+I$ to $I$. Note that there is no term constructor for the pure type $[s]\tau$, we must first bind the monadic computation $\store e$ to obtain the inner potential.

\begin{figure}[H]
\begin{mathpar}
 \inferrule*[right=T-Lam]
      {\hastype{\Theta;\Delta;\Omega;\Gamma, x:\tau_1}{e}{\tau_2}}
      {\hastype{\Theta;\Delta;\Omega;\Gamma}{\tlam{x}{e}}{\tau_1 \multimap \tau_2}}

\inferrule*[right=T-App]
      {\hastype{\Theta;\Delta;\Omega;\Gamma_1}{e_1}{\tau_1 \multimap \tau_2}
       \quad \hastype{\Theta;\Delta;\Omega;\Gamma_2}{e_2}{\tau_1}}
      {\hastype{\Theta;\Delta;\Omega;\Gamma_1+\Gamma_2}{e_1e_2}{\tau_2}}

\inferrule*[right=T-ExpI]
  {\hastype{\Theta;\Delta;\Omega;\cdot}{e}{\tau}}
  {\hastype{\Theta;\Delta;\Omega;\cdot}{!e}{!\tau}}

\inferrule*[right=T-ExpE]
  {\hastype{\Theta;\Delta;\Omega;\Gamma_1}{e_1}{!\tau_1}
    \quad\hastype{\Theta;\Delta;\Omega,x:\tau_1;\Gamma_2}{e_2}{\tau_2}}
  {\hastype{\Theta;\Delta;\Omega;\Gamma_1+\Gamma_2}
    {\lete{!x}{e_1}{e_2}}{\tau_2}}

\inferrule*[right=T-TensorI]
      {\hastype{\Theta;\Delta;\Omega;\Gamma_1}{e_1}{\tau_1}
       \quad \hastype{\Theta;\Delta;\Omega;\Gamma_2}{e_2}{\tau_2}}
      {\hastype{\Theta;\Delta;\Omega;\Gamma_1+\Gamma_2}{\ttuple{e_1}{e_2}}{\tau_1\otimes\tau_2}}

\inferrule*[right=T-TensorE]
  {\hastype{\Theta;\Delta;\Omega;\Gamma_1}{e_1}{\tau_1\otimes\tau_2}
    \quad\hastype{\Theta;\Delta;\Omega;\Gamma_2,x:\tau_1,y:\tau_2}{e_2}{\tau}}
  {\hastype{\Theta;\Delta;\Omega;\Gamma_1+\Gamma_2}
    {\lete{\ttuple{x}{y}}{e_1}{e_2}}{\tau}} 

\inferrule*[right=T-WithI]
      {\hastype{\Theta;\Delta;\Omega;\Gamma}{e_1}{\tau_1}
       \quad \hastype{\Theta;\Delta;\Omega;\Gamma}{e_2}{\tau_2}}
      {\hastype{\Theta;\Delta;\Omega;\Gamma}{\tuple{e_1}{e_2}}{\tau_1\;\&\;\tau_2}}

\inferrule*[right=T-Fst]
      {\hastype{\Theta;\Delta;\Omega;\Gamma}{e}{\tau_1\;\&\;\tau_2}}
      {\hastype{\Theta;\Delta;\Omega;\Gamma}{\mathsf{fst}\,e}{\tau_1}}

\inferrule*[right=T-Snd]
      {\hastype{\Theta;\Delta;\Omega;\Gamma}{e}{\tau_1\;\&\;\tau_2}}
      {\hastype{\Theta;\Delta;\Omega;\Gamma}{\mathsf{snd}\,e}{\tau_2}}

\inferrule*[right=T-Store]
  {\hastype{\Theta;\Delta;\Omega;\Gamma}{e}{\tau}
    \quad\Theta;\Delta\vdash s:\mathbb{R}_{\ge0}}
  {\hastype{\Theta;\Delta;\Omega;\Gamma}{\store e}{\tmonad{s}{([s]\tau)}}}

\inferrule*[right=T-Release]
  {\hastype{\Theta;\Delta;\Omega;\Gamma_1}{e_1}{[s]\tau_1}
    \quad\hastype{\Theta;\Delta;\Omega;\Gamma_2,x:\tau_1}{e_2}
      {\tmonad{(s+I)}{\tau_2}}}
  {\hastype{\Theta;\Delta;\Omega;\Gamma_1+\Gamma_2}
    {\release{x}{e_1}{e_2}}{\tmonad{I}{\tau_2}}}
\end{mathpar} 
\caption{Selected typing rules for $\lambda^g_{\mathsf{amor}}$}
\label{fig:amor-typing-rules}
\end{figure}

\subsubsection{Type Formations and Subtyping}
The well-formedness and subtyping rules of $\lambda^g_{\mathsf{amor}}$ largely follow those of $\lambda^g_{\mathsf{worst}}$ presented in Section~\ref{sec:worst-case-type-system},
with standard extensions for the new types added. 
The notable rules are type formation and subtyping for the potential type, presented in Figure \ref{fig:amor-wellform-sub}.

Intuitively, the potential index $s$ is required to be a real number and consistent with the cost grades of monads.
However, we also restrict it to be time-independent by the familiar constraint $s = \mathsf{S}(s)$.
The intuition of this is similar to that of affine typing: potential should not be created arbitrarily. 
For example, the index $\clk$ that carries $i$ units of potential at time $i$ will be regarded as carrying $i+1$ units one step later, 
without any additional potential being stored. Specifically, using time-dependent potential will lead to unsoundness of the type system, we wiil
show a corresponding program example in the next section.
\begin{figure}[H]
\begin{mathpar}
    \inferrule*[right=WF-Pot]
      {\Psi;\Theta;\Delta\vdash\tau \quad 
        \Theta;\Delta\vdash s:\mathbb{R}_{\ge0}
        \quad \Theta;\Delta \models s=\mathsf{S}(s) 
        }
      {\Psi;\Theta;\Delta\vdash[s]\tau}

    \inferrule*[right=Sub-Pot]
    {\Theta;\Delta\vdash\tau<:\tau'\quad\Theta;\Delta\models s'\leq s}
    {\Theta;\Delta\vdash[s]\tau<:[s']\tau'}
\end{mathpar}
\caption{Type Well-formedness and Subtyping rules for potential}
\label{fig:amor-wellform-sub}
\end{figure}
The \textsc{Sub-Pot} rule allows excess potential to be discarded: a term having $s$ units of potential can be treated as carrying any smaller amount $s'$. 
This is dual to \textsc{Sub-Monad}, which allows cost upper bounds to be increased.

\subsection{Semantics}
The operational semantics of $\lambda^g_{\mathsf{amor}}$ is akin to that of $\lambda^g_{\mathsf{worst}}$, 
with additional rules for the new constructs. Here, we focus on the forcing evaluation rules for terms related to the potential type, presented in Figure \ref{fig:amor-forcing}.

Although $\store e$ is statically assigned a cost accounting for the potential stored, its forcing evaluation incurs zero operational cost.
The $\mathsf{release}$ operation resembles $\mathsf{bind}$, which just sequences two computations with the cost of the second.
The semantics is sound since storing potential costs nothing, and hence releasing it to a computation does not decrease the cost.

This reveals the ``ghost'' nature of potential, where it has no runtime manifestation, and the type system ensures that 
it is used soundly.

\begin{figure}[H]
\begin{mathpar}
  \inferrule*[right=E-Store]
  {e\Downarrow_n v}
  {\store e\Downarrow_{n,i}^{0}v}

\inferrule*[right=E-Release]
  { e_1\Downarrow_n v_1 \quad
    e_2[v_1/x]\Downarrow_n v_2 
    \quad v_2\Downarrow_{n,i}^{c}v}
  {\release{x}{e_1}{e_2}\Downarrow_{n,i}^{c}v}
\end{mathpar}
\caption{Forcing evaluation rules for $\store$ and $\mathsf{release}$}
\label{fig:amor-forcing}
\end{figure}

An unsound example: 
$ \\
\mathsf{releaseClk}: \tmonad{\mathsf{clk}}{\textbf{1}} \\
\mathsf{releaseClk} \triangleq \bind{p}{\store-}{\pull}(
\delay (\release{-}{p}{\tick{\mathsf{clk}}})) \\
\mathsf{releaseClk} \Downarrow^{\mathsf{clk}(i+1)}_{n+1,i} \delay\langle\rangle  \\
$







